\vspace{-0.1cm}
\section{Preliminary}
\label{sec:prelim}
\vspace{-0.2cm}



% diffusion Diffusion 
% clip img encoder, text enc
% text embedding: 
% ENC_{img}(I), ENC_{text}(T)
% remove A



% \textbf{Functional Magnetic Resonance Imaging (fMRI). } 
% Each raw fMRI sample is a 4D volume (space and time), averaged over time to produce a 3D brain volume per stimulus. This is then flattened into a 1D vector to work with machine learning models, with each element representing a single voxel’s neural activation~\citep{allen2022massive}.



\textbf{Notations. } 
% \yj{We denote the set of fMRI samples collected while the subject views images as $\mathcal{X}$. Each sample is preprocessed following \citep{scotti2023reconstructing}, yielding a 1D vector $\mathbf{x}_i \in \mathbb{R}^v$ representing neural activity across $v$ voxels. We partition $\mathcal{X}$ into training, validation, and test subsets, using superscripts to indicate the split—for example, the training set with $N$ samples is denoted as $\mathcal{X}^{\text{train}} = \{\mathbf{x}_1, \ldots, \mathbf{x}_N\}$.
% % Let $\mathcal{X}^{\text{train}} = \{\mathbf{x}_1, \ldots, \mathbf{x}_N\}$ and $\mathcal{X}^{\text{test}} = \{\mathbf{x}_1, \ldots, \mathbf{x}_M\}$ denote the sets of training and test fMRI samples, respectively. We assume that the fMRI data is collected while the subject is viewing an image and is preprocessed following \citep{scotti2023reconstructing}, resulting in each fMRI sample being represented as a 1D vector $\mathbf{x}_i \in \mathbb{R}^v$ with $v$ voxels.
% We denote the $i$-th image stimulus presented during fMRI recording as $\mathbf{Y}_i \in \mathbb{R}^{H \times W \times 3}$. Let $m$ represent the number of objects in the image.
% % as identified by a pre-trained \yj{don't mention identified by a pre-trained model, as the number of objects should be something objective. you use vlm to detect it, but that is just a tool.}
% % \xiang{object recognition?} 
% % model. 
% % \yj{the next sentence is a little bit surprising. add a transition sentence, say that in your method, you will use text descriptions of the objects in the image, which you utilize a VLM to generate it. } 
% In our method, we use a pre-trained vision-language model (VLM) to generate object-level descriptions from each image. An attribute keyword $a$—a text phrase specifying aspects such as color, geometry, or layout—is used to construct a prompt $\mathcal{P}(a)$. Given an image $\mathbf{Y}_i$ and prompt $\mathcal{P}(a)$, the VLM produces descriptions conditioned on the specified attribute.}
% %The resulting structured output is denoted by $D_i^a=\text{VLM}(I_i,\mathcal{P}(a))$.
In our work, we denote the set of fMRI samples collected during image viewing as $\mathcal{X}$. Each sample is a preprocessed 1D vector $\mathbf{x}_i \in \mathbb{R}^v$, capturing neural activity across $v$ voxels selected from brain regions \citep{scotti2023reconstructing}. The dataset is split into training and test subsets, with superscripts indicating the split. For example, we denote the training set with $N$ samples as: $\mathcal{X}^{\text{train}} = \{\mathbf{x}_1, \ldots, \mathbf{x}_N\}$. The corresponding image stimulus for the $i$-th sample is denoted as: $\mathbf{Y}_i \in \mathbb{R}^{H \times W \times 3}$, which contains $m$ objects. 
% To generate object-level textual descriptions of each image, we use a pre-trained vision-language model (VLM) guided by an attribute keyword $a$, which specifies a property such as geometry, action, or spatial layout. We construct a prompt $\mathcal{P}(a)$ from $a$ and feed it along with $\mathbf{Y}_i$ into the VLM to generate attribute-conditioned descriptions.

%which is a hierarchical textual description that includes the most salient objects in the image and provides detailed, attribute-conditioned descriptions for each. Specifically, it consists of a list of object-level entries, where each entry includes a natural language description annotated with semantic roles (e.g., [Agent], [Action], [Patient]). These role labels are derived from PropBank-style semantic analysis and serve to disentangle the compositional meaning of each object’s description.
%We use $h_i^{I}, h_i^{T}$ to denote the flattened representation extract from model $m$. The model could be Clip image encoder,  Clip text encoder, Resnet50 

% Let $\mathcal{A} = \{a_1, a_2, \dots, a_p\}$ denote a set of image attributes, where each attribute $a_i$ can be used to describe an image $\mathbf{G}_n$. Given $\mathbf{G}_n$ and an attribute $a_i$, a vision-language model generates a corresponding textual description $d_i$, forming the set $\mathcal{D} = \{d_1, d_2, \dots, d_p\}$. We use $\mathbf{X}_{n} \in \mathbb{R}^d$ denote the preprocessed fMRI data recorded when a subject is viewing the $n$-th image $\mathbf{G}_n \in \mathbb{R}^{H \times W \times 3}$, where $d$ is the number of voxels, and $H$, $W$ are the height and width of the image, respectively. We use $\mathbf{F}_{n}^v$ as the $n$-th image representation extracted from vision model $v$ (e.g., CLIP Image encoder \citep{radford2021learning} , ResNet50 \citep{he2016deep}), and $\mathbf{F}_{n}^t$ as the $n$-th text representation extracted from language model $t$ (e.g., T5 \citep{raffel2020exploring}, LLama \citep{grattafiori2024llama} , CLIP Text encoder \citep{radford2021learning}). 

$\textbf{Problem Setup. }$ Our goal is to reconstruct the visual images that subjects viewed during fMRI recording. Formally, we seek to learn a reconstruction function $\mathcal{F}: \mathbb{R}^v \rightarrow \mathbb{R}^{H \times W \times 3}$ that maps each fMRI sample $\mathbf{x}_i$ to its corresponding image stimulus $\mathbf{Y}_i$.
% In our framework, we instantiate $\mathcal{F}$ as a composition of three components: an multi-layer perceptron (MLP) that maps $X_i$ into a language embedding space, a language model (LM) that generates structured descriptions, and a diffusion model that reconstructs the image:
% $
% \mathcal{F}(X_i) = \text{Diffusion}(\text{LM}(\text{MLP}(X_i))).
% $ The objective of the fmri-image reconstuction is to minimize the perceptual loss between $\hat{I}_i$ and $I_i$.

% diffusion pipeline, basic diffusion, various diffusion
$\textbf{Diffusion Model. }$ 
Diffusion models~\citep{rombach2022high, zhang2023adding} are a class of generative models that synthesize data by learning to reverse a multi-step noising process. Starting from Gaussian noise, they iteratively denoise a latent variable over a fixed number of timesteps using a denoising network—typically a U-Net—conditioned on the current timestep $t$. This process gradually produces samples resembling the training distribution.
 % \yj{lack citation for this conditional generation}
To incorporate external inputs such as text, the denoising network can take a conditioning input $\mathbf{C}$, usually text embeddings from a pre-trained encoder~\citep{zhang2023adding}. Conditioning is implemented via cross-attention mechanisms within the U-Net architecture~\citep{williams2023unified}, enabling the integration of textual information. In these layers, the latent representations $\mathbf{H}_t \in \mathbb{R}^{h \times w \times d}$ at time $t$ serve as queries, and $\mathbf{C} \in \mathbb{R}^{d_t \times d'}$ serves as both keys and values~\citep{williams2023unified, yang2024mastering}: 
%\yj{check formula and meaning}
\begin{equation*}
\vspace{-0.1cm}
\text{CrossAttention}(\mathbf{H}_t, \mathbf{C}) = \text{softmax} \left( \frac{ \phi(\mathbf{H}_t) \cdot \mathbf{W}_Q \cdot (\varphi(\mathbf{C}) \cdot \mathbf{W}_{K})^\top}{\sqrt{d_k}} \right) \varphi(\mathbf{C}) \cdot \mathbf{W}_{V},
\end{equation*}
where $\mathbf{W}_Q\in \mathbb{R}^{d \times d_k} $ ,  $\mathbf{W}_K \in \mathbb{R}^{d^\prime \times d_k} $, $\mathbf{W}_V \in \mathbb{R}^{d^\prime \times d}$ are projection matrices; and $ \phi(\cdot)$ and $ \varphi(\cdot)$ are learned transformations.

